% problem_id: S001-lemma-sum-is-ok
% source_id: S001 (Stacks Project)
% source_file: more-morphisms.tex
% statement_locator: more-morphisms.tex lines 8595--8610
% proof_locator: more-morphisms.tex lines 8612--8777
% env: lemma
% id_note: label 在本来源内唯一
% pairing: tight (verified)
% status: complete (statement + author proof verbatim + reason + locators)
%
\begin{lemma}
\label{lemma-sum-is-ok}
Let $R \to S$ be a finite type ring map. Let $d \geq 0$. Let $a, b \in S$.
Assume that the fibres of
$$
f_a : \Spec(S) \longrightarrow \mathbf{A}^1_R
$$
given by the $R$-algebra map $R[x] \to S$ sending $x$ to $a$
have dimension $\leq d$. Then there exists an $n_0$ such that for
$n \geq n_0$ the fibres of
$$
f_{a^n + b} : \Spec(S) \longrightarrow \mathbf{A}^1_R
$$
given by the $R$-algebra map $R[x] \to S$ sending $x$ to $a^n + b$
have dimension $\leq d$.
\end{lemma}

\begin{proof}
In this paragraph we reduce to the case where $R \to S$ is of finite
presentation. Namely, write $S = R[A, B, x_1, \ldots, x_n]/J$ for some ideal
$J \subset R[x_1, \ldots, x_n]$ where $A$ and $B$ map to $a$ and $b$ in $S$.
Then $J$ is the union of its finitely generated ideals $J_\lambda \subset J$.
Set $S_\lambda = R[A, B, x_1, \ldots, x_n]/J_\lambda$ and
denote $a_\lambda, b_\lambda \in S_\lambda$ the images of $A$ and $B$.
Then for some $\lambda$ the fibres of
$$
f_{a_\lambda} : \Spec(S_\lambda) \longrightarrow \mathbf{A}^1_R
$$
have dimension $\leq d$, see Limits, Lemma \ref{limits-lemma-limit-dimension}.
Fix such a $\lambda$. If we can find $n_0$ which works for $R \to S_\lambda$,
$a_\lambda$, $b_\lambda$, then $n_0$ works for $R \to S$. Namely, the fibres
of $f_{a_\lambda^n + b_\lambda} : \Spec(S_\lambda) \to \mathbf{A}^1_R$
contain the fibres of $f_{a^n + b} : \Spec(S) \to \mathbf{A}^1_R$.
This reduces us to the case discussed in the next paragraph.

\medskip\noindent
Assume $R \to S$ is of finite presentation. In this paragraph we reduce
to the case where $R$ is of finite type over $\mathbf{Z}$. By
Algebra, Lemma \ref{algebra-lemma-limit-module-finite-presentation}
we can find a directed set $\Lambda$ and a system of ring maps
$R_\lambda \to S_\lambda$ over $\Lambda$ whose colimit is
$R \to S$ such that $S_\mu = S_\lambda \otimes_{R_\lambda} R_\mu$
for $\mu \geq \lambda$ and such that each $R_\lambda$ and $S_\lambda$
is of finite type over $\mathbf{Z}$. Choose $\lambda_0 \in \Lambda$ and
elements $a_{\lambda_0}, b_{\lambda_0} \in S_{\lambda_0}$ mapping to
$a, b \in S$. For $\lambda \geq \lambda_0$ denote
$a_\lambda, b_\lambda \in S_\lambda$
the image of $a_{\lambda_0}, b_{\lambda_0}$.
Then for $\lambda \geq \lambda_0$ large enough the fibres of
$$
f_{a_\lambda} : 
\Spec(S_\lambda) \longrightarrow \mathbf{A}^1_{R_\lambda}
$$
have dimension $\leq d$, see
Limits, Lemma \ref{limits-lemma-descend-dimension-d}.
Fix such a $\lambda$. If we can find $n_0$ which works for
$R_\lambda \to S_\lambda$, $a_\lambda$, $b_\lambda$, then $n_0$
works for $R \to S$. Namely, any fibre of
$f_{a^n + b} : \Spec(S) \to \mathbf{A}^1_R$
has the same dimension as a fibre of
$f_{a_\lambda^n + b_\lambda} : \Spec(S_\lambda) \to \mathbf{A}^1_{R_\lambda}$
by Morphisms, Lemma \ref{morphisms-lemma-dimension-fibre-after-base-change}.
This reduces us the the case discussed in the next paragraph.

\medskip\noindent
Assume $R$ and $S$ are of finite type over $\mathbf{Z}$. In particular
the dimension of $R$ is finite, and we may use induction on $\dim(R)$.
Thus we may assume the result holds for all situations with
$R' \to S'$, $a$, $b$ as in the lemma with $R'$ and $S'$ of finite type
over $\mathbf{Z}$ but with $\dim(R') < \dim(R)$.

\medskip\noindent
Since the statement is about the topology of the spectrum of $S$
we may assume $S$ is reduced. Let $S^\nu$ be the normalization of $S$.
Then $S \subset S^\nu$ is a finite extension as $S$ is excellent, see
Algebra, Proposition \ref{algebra-proposition-ubiquity-nagata}
and Morphisms, Lemma \ref{morphisms-lemma-nagata-normalization}.
Thus $\Spec(S^\nu) \to \Spec(S)$ is surjective and finite
(Algebra, Lemma \ref{algebra-lemma-integral-overring-surjective}).
It follows that if the result holds for $R \to S^\nu$ and the
images of $a$, $b$ in $S^\nu$, then the result holds for $R \to S$, $a$, $b$.
(Small detail omitted.) This reduces us to the case discussed in the
next paragraph.

\medskip\noindent
Assume $R$ and $S$ are of finite type over $\mathbf{Z}$ and
$S$ normal. Then $S = S_1 \times \ldots \times S_r$ for some
normal domains $S_i$. If the result holds for each $R \to S_i$
and the images of $a$, $b$ in $S_i$, then the result holds for
$R \to S$, $a$, $b$. (Small detail omitted.)
This reduces us to the case discussed in the next paragraph.

\medskip\noindent
Assume $R$ and $S$ are of finite type over $\mathbf{Z}$ and
$S$ a normal domain. We may replace $R$ by the image of $R$ in $S$
(this does not increase the dimension of $R$).
This reduces us to the case discussed in the next paragraph.

\medskip\noindent
Assume $R \subset S$ are of finite type over $\mathbf{Z}$ and
$S$ a normal domain. Consider the morphism
$$
f_a : \Spec(S) \to \mathbf{A}^1_R
$$
The assumption tells us that $f_a$ has fibres of dimension $\leq d$.
Hence the fibres of $f : \Spec(S) \to \Spec(R)$ have dimension $\leq d + 1$
(Morphisms, Lemma \ref{morphisms-lemma-dimension-fibre-at-a-point-additive}).
Consider the morphism of integral schemes
$$
\phi : \Spec(S) \to \mathbf{A}^2_R = \Spec(R[x, y])
$$
corresponding to the $R$-algebra map $R[x, y] \to S$ sending $x$ to $a$
and $y$ to $b$. There are two cases to consider
\begin{enumerate}
\item $\phi$ is dominant, and
\item $\phi$ is not dominant.
\end{enumerate}
We claim that in both cases there exists an integer $n_0$
and a nonempty open $V \subset \Spec(R)$ such that for $n \geq n_0$
the fibres of $f_{a^n + b}$ at points $q \in \mathbf{A}^1_V$
have dimension $\leq d$.

\medskip\noindent
Proof of the claim in case (1). We have $f_{a^n + b} = \pi_n \circ \phi$
where
$$
\pi_n :  \mathbf{A}^2_R \to \mathbf{A}^1_R
$$
is the flat morphism corresponding to the $R$-algebra map $R[x] \to R[x, y]$
sending $x$ to $x^n + y$. Since $\phi$ is dominant there is a dense open
$U \subset \Spec(S)$ such that $\phi|_U : U \to \mathbf{A}^2_R$ is flat
(this follows for example from generic flatness, see
Morphisms, Proposition \ref{morphisms-proposition-generic-flatness}).
Then the composition
$$
f_{a^n + b}|_U :
U \xrightarrow{\phi|_U}
\mathbf{A}^2_R
\xrightarrow{\pi_n}
\mathbf{A}^1_R
$$
is flat as well. Hence the fibres of this morphism have at least
codimension $1$ in the fibres of $f|_U : U \to \Spec(R)$ by
Morphisms, Lemma \ref{morphisms-lemma-dimension-fibre-at-a-point-additive}.
In other words, the fibres of $f_{a^n + b}|_U$ have dimension $\leq d$.
On the other hand, since $U$ is dense in $\Spec(S)$, we can find a nonempty
open $V \subset \Spec(R)$ such that $U \cap f^{-1}(p) \subset f^{-1}(p)$
is dense for all $p \in V$ (see for example Lemma
\ref{lemma-nowhere-dense-generic-fibre}).
Thus $\dim(f^{-1}(p) \setminus U \cap f^{-1}(p)) \leq d$ and
we conclude that our claim is true (as any fibres of
$f_{a^n + b} : \Spec(S) \to \mathbf{A}^1_R$
is contained in a fibre of $f$).

\medskip\noindent
Case (2). In this case we can find a nonzero $g = \sum c_{ij} x^i y^j$
in $R[x, y]$ such that $\Im(\phi) \subset V(g)$. In fact, we may assume
$g$ is irreducible over $\text{Frac}(R)$. If $g \in R[x]$, say with leading
coefficient $c$, then over $V = D(c) \subset \Spec(R)$ the fibres of
$f$ already have dimension $\leq d$ (because the image of $f_a$ is contained
in $V(g) \subset \mathbf{A}^1_R$ which has finite fibres over $V$). Hence
we may assume $g$ is not contained in $R[x]$.
Let $s \geq 1$ be the degree of $g$ as a polynomial in $y$ and
let $t$ be the degree of $\sum c_{is} x^i$ as a polynomial in $x$.
Then $c_{ts}$ is nonzero and
$$
g(x, -x^n) = (-1)^s c_{ts} x^{t + sn} + l.o.t.
$$
provided that $n$ is bigger than the degree of $g$ as a polynomial in $x$
(small detail omitted). For such $n$ the polynomial $g(x, -x^n)$ 
is a nonzero polynomial in $x$ and maps to a nonzero polynomial in
$\kappa(\mathfrak p)[x]$ for all $\mathfrak p \subset R$,
$c_{st} \not \in \mathfrak p$. We conclude that our claim is
true for $V$ equal to the principal open $D(c_{ts})$ of $\Spec(R)$.

\medskip\noindent
OK, and now we can use induction on $\dim(R)$. Namely, let
$I \subset R$ be an ideal such that $V(I) = \Spec(R) \setminus V$.
Observe that $\dim(R/I) < \dim(R)$ as $R$ is a domain.
Let $n'_0$ be the integer we have by induction on $\dim(R)$
for $R/I \to S/IS$ and the images of $a$ and $b$ in $S/IS$.
Then $\max(n_0, n'_0)$ works.
\end{proof}
